\documentclass[11pt,a4paper]{article}
\usepackage[utf8]{inputenc}
\usepackage[T1]{fontenc}
\usepackage{amsmath,amssymb,amsthm}
\usepackage{graphicx}
\usepackage{hyperref}
\usepackage{geometry}
\usepackage{booktabs}
\usepackage{enumitem}

\geometry{margin=2.5cm}

\newtheorem{theorem}{Theorem}
\newtheorem{lemma}[theorem]{Lemma}
\newtheorem{proposition}[theorem]{Proposition}
\newtheorem{corollary}[theorem]{Corollary}
\theoremstyle{definition}
\newtheorem{definition}[theorem]{Definition}
\theoremstyle{remark}
\newtheorem{remark}[theorem]{Remark}

\title{\textbf{The Geometric Origin of Residue-Class Gap Asymmetry\\in Sparse Arithmetic Point Processes}}

\author{Thomas Hoffbauer\\
\small \texttt{(Author affiliation)}}

\date{June 23, 2026}

\begin{document}

\maketitle

\begin{abstract}
We study the conditional gap distribution $P(g \bmod 12 \mid p_n \equiv a)$ for consecutive primes in residue classes modulo 12 and observe a systematic asymmetry: $P(g \equiv 2,4 \mid a) > P(g \equiv 8,10 \mid a)$ for all four non-trivial residue classes $a \in \{1, 5, 7, 11\}$.

\textbf{We identify the residue-class gap asymmetry as a universal consequence of geometric gap statistics rather than a uniquely prime-specific phenomenon.} For Bernoulli point processes (discrete Poisson analogues) with constant acceptance probability $p$, we derive the exact formula $R = (1-p)^{-6}$, where $R$ is the ratio of gap probabilities. Our data and model analyses suggest that this geometric origin accounts for a substantial portion of the observed asymmetry in primes.

Empirical comparison with three null models reveals an unexpected hierarchy: $R_{\text{Bernoulli}} \approx 1.83 > R_{\text{Prime}} = 1.58 > R_{\text{Cramér}} \approx 1.36$. This indicates that the naive Cramér model (logarithmic acceptance $p(n) = 2/\ln(n)$) \emph{overdamps} the asymmetry through mixing, while real primes show intermediate behavior consistent with sieve-induced correlations.

Our results place primes in an intermediate regime between uncorrelated Bernoulli processes (maximal short-gap bias) and the overdamped logarithmic-density dynamics of the Cramér model, suggesting that arithmetic correlations (sieve structure) partially restore the asymmetry suppressed by logarithmic mixing.

\textbf{Keywords:} Prime gaps, residue classes, Bernoulli process, Cramér model, gap asymmetry, geometric distribution
\end{abstract}

\section{Introduction}

\subsection{Motivation}

The distribution of prime numbers has fascinated mathematicians for millennia. While the Prime Number Theorem establishes the asymptotic density $\pi(x) \sim x/\ln(x)$, the \emph{local} structure of consecutive primes — their gaps $g_n = p_{n+1} - p_n$ — remains rich with patterns and open questions.

Recent investigations \cite{Rubinstein1994,Granville2006} have revealed systematic biases in prime gap distributions, notably the Chebyshev bias and prime number races. These phenomena show that primes, despite having equal limiting density in residue classes, exhibit \emph{persistent statistical asymmetries} over finite ranges.

\subsection{The Observation}

We classify primes $p > 3$ by their residue modulo 12:
\begin{align*}
E: p \equiv 1 \pmod{12}, \quad
A: p \equiv 5 \pmod{12}, \quad
B: p \equiv 7 \pmod{12}, \quad
C: p \equiv 11 \pmod{12}
\end{align*}

For consecutive primes $(p_n, p_{n+1})$ with gap $g = p_{n+1} - p_n$, we observe a \textbf{systematic asymmetry} in the conditional gap distribution:

\begin{equation}
P(g \equiv 2,4 \bmod 12 \mid p_n \equiv a) > P(g \equiv 8,10 \bmod 12 \mid p_n \equiv a)
\end{equation}

for \emph{all} starting classes $a \in \{E, A, B, C\}$.

\subsection{The Question}

Is this asymmetry:
\begin{enumerate}[label=(\alph*)]
\item A deep property specific to prime numbers (sieve structure, Hardy-Littlewood correlations)?
\item A general phenomenon of sparse point sets (sparsity-induced)?
\item Something in between?
\end{enumerate}

\subsection{Our Answer}

\textbf{Answer: (c) — but with a surprising twist.}

We show that:
\begin{itemize}
\item The asymmetry \emph{already exists} in Bernoulli point processes (discrete Poisson analogues) with geometric gap distributions
\item It is \emph{not} prime-specific in the strong sense
\item But the \emph{magnitude} depends critically on the acceptance model
\item Primes occupy an \textbf{intermediate regime} between uncorrelated Bernoulli processes and overdamped Cramér dynamics
\end{itemize}

\subsection{Main Results}

\begin{theorem}[Geometric Asymmetry Formula]
Let $\mathcal{B}$ be a Bernoulli point process (discrete Poisson analogue) on $\mathbb{N}$ with constant acceptance probability $p$. Define the gap asymmetry ratio:
\begin{equation}
R_{\mathcal{B}} = \frac{P(g \equiv 2 \text{ or } 4 \bmod 12)}{P(g \equiv 8 \text{ or } 10 \bmod 12)}
\end{equation}
Then:
\begin{equation}\label{eq:main_formula}
\boxed{R_{\mathcal{B}} = (1-p)^{-6}}
\end{equation}
\end{theorem}

\begin{theorem}[Null Model Hierarchy]
For three sparse point processes up to $10^5$ with identical modulo-12 evaluation:
\begin{equation}
R_{\text{Bernoulli}} \approx 1.83 \quad > \quad R_{\text{Prime}} = 1.58 \quad > \quad R_{\text{Cramér}} \approx 1.36
\end{equation}
This hierarchy is \textbf{stable over 10 random seeds} (std.~dev.~$< 0.07$).
\end{theorem}

\subsection{Interpretation}

The hierarchy reveals a \textbf{two-stage effect}:

\begin{enumerate}
\item \textbf{Bernoulli $\to$ Cramér ($-28\%$):} Logarithmic acceptance $p(n) = 2/\ln(n)$ creates \emph{mixing} over variable density, strongly damping the geometric short-gap bias.

\item \textbf{Cramér $\to$ Prime ($+16\%$):} Primes exhibit intermediate behavior. This is consistent with sieve structure (twins, cousins, sexy primes) enhancing short gaps relative to independent acceptance, though quantitative derivation remains open (Section 7.3).
\end{enumerate}

\textbf{Conclusion:} Primes lie in an intermediate regime between uncorrelated Bernoulli processes (maximal geometric bias) and the overdamped logarithmic-density dynamics of Cramér.

\section{The Chirality Observable}

\subsection{Complete EABC Quadruples}

\begin{definition}
A prime quadruple $Q = (p, q, r, s)$ is \textbf{complete} if:
\begin{equation}
\{k(p), k(q), k(r), k(s)\} = \{E, A, B, C\}
\end{equation}
where $k(p)$ denotes the EABC-class of prime $p$.
\end{definition}

For consecutive complete quadruples $Q_n = (p_n, p_{n+1}, p_{n+2}, p_{n+3})$, we define the \textbf{chirality observable}:

\begin{equation}
\chi(Q) = \begin{cases}
+1, & \text{if norm}(\text{sig}(Q)) = EABC \\
-1, & \text{if norm}(\text{sig}(Q)) = ECBA \\
0, & \text{otherwise}
\end{cases}
\end{equation}

where norm$(\cdot)$ cyclically rotates the signature to start with $E$.

\subsection{Orientation Bias}

The \textbf{bias function} is:
\begin{equation}
H(X) = \sum_{Q \leq X} \chi(Q) = N_{\text{EABC}}(X) - N_{\text{ECBA}}(X)
\end{equation}

\textbf{Empirical observation (up to $X = 10^5$):}
\begin{equation}
h(X) = \frac{H(X)}{N(X)} \approx 0.16 \quad \text{(persistent bias)}
\end{equation}

\subsection{Connection to Gap Asymmetry}

The key insight: The orientation bias \emph{follows mechanistically} from gap asymmetry via the deterministic relation:
\begin{equation}
p_{n+1} \equiv p_n + g_n \pmod{12}
\end{equation}

\textbf{EABC cycle} uses gaps $\{4, 2, 4, 2\} \subset \{2, 4\}$.\\
\textbf{ECBA cycle} uses gaps $\{10, 8, 10, 8\} \subset \{8, 10\}$.

Therefore:
\begin{equation}
\frac{P(\text{EABC})}{P(\text{ECBA})} \approx \left(\frac{P(g \equiv 2,4)}{P(g \equiv 8,10)}\right)^4 \approx R^4
\end{equation}

Thus, the \textbf{chirality observable is a probe} for the underlying gap asymmetry.

\section{Bernoulli Point Process Model: The Geometric Origin}

\subsection{Setup}

Consider a Bernoulli point process (discrete Poisson analogue) on $\mathbb{N}$ where each candidate point is accepted with constant probability $p$.

\textbf{Gap distribution:} If a point is accepted at position $n$, the gap $G$ to the next accepted point is \textbf{geometrically distributed}:
\begin{equation}
P(G = k) = p(1-p)^{k-1} = p q^{k-1}, \quad q := 1-p
\end{equation}

\subsection{Residue Class Distribution}

For gaps modulo 12, we have:
\begin{equation}
P(G \equiv r \bmod 12) = \sum_{\substack{k \geq 1 \\ k \equiv r \pmod{12}}} p q^{k-1}
= p q^{r-1} \sum_{j=0}^{\infty} q^{12j}
= \frac{p q^{r-1}}{1 - q^{12}}
\end{equation}

\textbf{Key observation:} The probability \emph{decreases exponentially} with residue class $r$!

\subsection{The Asymmetry Ratio}

For our specific residue classes:
\begin{align}
P(G \equiv 2 \bmod 12) &= \frac{p q^{1}}{1-q^{12}} \\
P(G \equiv 4 \bmod 12) &= \frac{p q^{3}}{1-q^{12}} \\
P(G \equiv 8 \bmod 12) &= \frac{p q^{7}}{1-q^{12}} \\
P(G \equiv 10 \bmod 12) &= \frac{p q^{9}}{1-q^{12}}
\end{align}

Therefore:
\begin{align}
R &= \frac{P(G \equiv 2) + P(G \equiv 4)}{P(G \equiv 8) + P(G \equiv 10)} \\
&= \frac{p(q^1 + q^3)/(1-q^{12})}{p(q^7 + q^9)/(1-q^{12})} \\
&= \frac{q^1(1 + q^2)}{q^7(1 + q^2)} \\
&= q^{-6}
\end{align}

\begin{equation}
\boxed{R_{\text{Bernoulli}} = (1-p)^{-6}}
\end{equation}

\subsection{Interpretation}

The asymmetry arises from \textbf{short-gap bias}: Geometric distributions favor small gaps, and residue classes $\{2, 4\}$ are \emph{systematically closer to 0} than $\{8, 10\}$ modulo 12.

This is \textbf{universal} for any constant-$p$ sparse process with geometric gaps.

\subsection{Numerical Prediction}

For our empirical Bernoulli processes (mean gap $\approx \ln(50000) \approx 10.82$):
\begin{equation}
p \approx \frac{1}{10.82} \approx 0.092
\end{equation}

Predicted:
\begin{equation}
R_{\text{Bernoulli}} = (1 - 0.092)^{-6} \approx 1.85
\end{equation}

\textbf{Observed:} $R_{\text{Bernoulli}} \approx 1.83$ (mean over 10 seeds)

\textbf{Agreement: Excellent!}

\section{Cramér Model: Logarithmic Damping}

\subsection{The Cramér Construction}

The naive Cramér model accepts each odd candidate $n$ independently with probability:
\begin{equation}
p(n) = \frac{2}{\ln(n)}
\end{equation}

This gives correct asymptotic density $\sim x/\ln(x)$ but \emph{lacks} sieve structure.

\subsection{Why Is It Too Tame?}

\textbf{Key difference:} $p(n)$ is \emph{not constant} but \emph{decreases with $n$}.

For $n \in [5, 100000]$:
\begin{itemize}
\item At $n=1000$: $p \approx 0.289$
\item At $n=50000$: $p \approx 0.185$
\item At $n=100000$: $p \approx 0.173$
\end{itemize}

\textbf{Naive prediction:} Mean $\bar{p} \approx 0.20$ would give:
\begin{equation}
R_{\text{naive}} = (1 - 0.20)^{-6} \approx 3.81
\end{equation}

\textbf{But we observe:} $R_{\text{Cramér}} \approx 1.36$

\textbf{Discrepancy:} Factor of 2.8!

\subsection{The Mixing Effect}

The gap distribution is \emph{not} purely geometric but a \textbf{mixture}:
\begin{equation}
P(G = k) = \int P(G = k \mid \text{start at } n) \, dP(n)
\end{equation}

Each window $[n, n+\Delta n]$ has different $p(n)$. This \textbf{averaging over variable acceptance rates} strongly damps extremes.

\textbf{Result:} The Cramér model \emph{overdamps} the asymmetry by mixing over variable density.

\section{Prime Data: The Empirical Hierarchy}

\subsection{Methodology}

We compare three models up to limit $10^5$:

\begin{enumerate}
\item \textbf{Bernoulli Process:} Geometric gap distribution with mean $\approx \ln(n)$ (constant $p$)
\item \textbf{Cramér Model:} Independent acceptance $p(n) = 2/\ln(n)$ (variable density)
\item \textbf{Real Primes:} Sieve of Eratosthenes
\end{enumerate}

All three undergo \emph{identical} modulo-12 gap analysis: Count transitions and compute:
\begin{equation}
R_{\text{gap}} = \text{geom.~mean}\left(\frac{P(g \equiv 2,4 \mid a)}{P(g \equiv 8,10 \mid a)}\right)
\end{equation}
over all four starting classes $a \in \{E, A, B, C\}$.

\subsection{Results}

\begin{table}[h]
\centering
\begin{tabular}{lccc}
\toprule
\textbf{Model} & \textbf{Sample Size} & \textbf{$R_{\text{gap}}$} & \textbf{Interpretation} \\
\midrule
Bernoulli (mean) & $\sim$6400 & 1.83 $\pm$ 0.06 & Geometric baseline \\
Real Primes & 9589 & 1.58 (fixed) & Intermediate regime \\
Cramér (mean) & $\sim$3200 & 1.36 $\pm$ 0.07 & Logarithmic overdamping \\
\bottomrule
\end{tabular}
\caption{Null model hierarchy over 10 random seeds (Bernoulli and Cramér) at $N=100{,}000$.}
\end{table}

\textbf{Hierarchy (consistent over all 10 seeds):}
\begin{equation}
\boxed{R_{\text{Bernoulli}} > R_{\text{Prime}} > R_{\text{Cramér}}}
\end{equation}

\textbf{Hierarchy (consistent over all 10 seeds):}
\begin{equation}
\boxed{R_{\text{Poisson}} > R_{\text{Prime}} > R_{\text{Cramér}}}
\end{equation}

\subsection{Quantitative Breakdown}

\begin{itemize}
\item \textbf{Poisson $\to$ Cramér:} $\Delta R = -0.47$ ($-28\%$) — Mixing overdamps
\item \textbf{Cramér $\to$ Prime:} $\Delta R = +0.22$ ($+16\%$) — Sieve restores
\item \textbf{Poisson $\to$ Prime:} $\Delta R = -0.25$ ($-14\%$) — Net regularization
\end{itemize}

\section{Interpretation: Structured Disorder}

\subsection{The Three Effects}

The hierarchy reveals \textbf{three distinct mechanisms}:

\begin{enumerate}
\item \textbf{Short-Gap Bias (Geometric Origin)}
\begin{itemize}
\item Source: Geometric/exponential distributions favor small gaps
\item Formula: $R = (1-p)^{-6}$ for constant $p$
\item Observable in: Pure Poisson process
\end{itemize}

\item \textbf{Logarithmic Thinning (Mixing Effect)}
\begin{itemize}
\item Source: $p(n) \sim 1/\ln(n)$ varies with $n$ $\to$ mixing
\item Effect: Dampens asymmetry via averaging
\item Observable in: Cramér model
\end{itemize}

\item \textbf{Arithmetic Back-Correlation (Sieve Effect)}
\begin{itemize}
\item Source: Hardy-Littlewood correlations, twin/cousin/sexy prime clustering
\item Effect: Enhances short gaps vs.~naive independent acceptance
\item Observable in: Real primes
\end{itemize}
\end{enumerate}

\subsection{The Equilibrium}

Primes sit at an \textbf{equilibrium} between:
\begin{align*}
\text{Geometric short-gap bias} &\quad (\text{from sparsity}) \\
\text{Logarithmic mixing} &\quad (\text{from density growth}) \\
\text{Sieve back-correlation} &\quad (\text{from arithmetic structure})
\end{align*}

This creates \textbf{structured disorder}: Not maximally random (vs.~Poisson), not smoothly damped (vs.~Cramér), but a \emph{middle regime} where arithmetic correlations tune the balance.

\subsection{Conceptual Spectrum}

\begin{equation}
\underbrace{\text{Bernoulli}}_{R \approx 1.83}
\quad \xrightarrow{\text{log-mixing}} \quad
\underbrace{\text{Cramér}}_{R \approx 1.36}
\quad \xrightarrow{\text{sieve}} \quad
\underbrace{\text{Primes}}_{R = 1.58}
\end{equation}

Primes are \textbf{more ordered than Bernoulli} (logarithmic growth regularizes) but \textbf{more disordered than Cramér} (sieve creates local clustering).

\section{Outlook and Open Questions}

\subsection{Hardy-Littlewood Quantification (Open Question)}

\textbf{The critical open question:} Quantify which portion of
\begin{equation}
\Delta R := R_{\text{Prime}} - R_{\text{Cramér}} \approx 0.22
\end{equation}
can be explained by Hardy-Littlewood $k$-tuple correlations.

\textbf{Proposed approach:} 
\begin{enumerate}
\item Compute the expected enhancement of short gaps (2, 4, 6) using:
\begin{itemize}
\item Twin prime constant: $\Pi_2 \approx 0.660$
\item Cousin prime constant: $\Pi_4$
\item Sexy prime constant: $\Pi_6$
\end{itemize}
\item Compare this theoretical prediction to the empirical gap $\Delta R$
\item If the Hardy-Littlewood correction quantitatively accounts for $\Delta R$, the causal chain
\begin{equation}
\text{Bernoulli} \xrightarrow{\text{log-mixing}} \text{Cramér} \xrightarrow{\text{sieve}} \text{Prime}
\end{equation}
would be established.
\end{enumerate}

\textbf{Current status:} This remains an open problem. The present work establishes the geometric baseline and empirical hierarchy, but does not yet close the quantitative gap.

\subsection{Modulo-30 Generalization (High Priority)}

The formula $R = (1-p)^{-\Delta r}$ should generalize to other moduli. For $m=30$ (8 residue classes), the predicted hierarchy should persist with adjusted exponents. \textbf{Testing this would significantly strengthen the universality claim.}

\subsection{Modified Cramér with Sieve}

\textbf{Prediction:} A modified Cramér model that \emph{explicitly excludes multiples} of small primes (sieve-Cramér hybrid) should satisfy:
\begin{equation}
R_{\text{Cramér+Sieve}} > R_{\text{Cramér}}
\end{equation}
approaching $R_{\text{Prime}}$.

\subsection{Asymptotic Behavior}

\textbf{Open Question:} Does $R(X)$ converge to a constant, or exhibit race-like oscillations as $X \to \infty$?

Current data show:
\begin{equation}
R_{\text{Prime}}(10^4) = 6.14 \to R_{\text{Prime}}(5 \times 10^5) = 3.59
\end{equation}

This suggests slow convergence, but the asymptotic limit remains open.

\section{Conclusion}

We identify the observed residue-class gap asymmetry $P(g \equiv 2,4) > P(g \equiv 8,10)$ in consecutive primes modulo 12 as arising substantially from the geometric distribution of gaps in sparse point processes, quantified by the exact formula $R = (1-p)^{-6}$ for Bernoulli point processes with constant acceptance probability.

Our data and model analyses indicate that the naive Cramér model \emph{overdamps} this asymmetry through mixing over variable acceptance probability $p(n) = 2/\ln(n)$, while real primes exhibit intermediate behavior. This places primes between uncorrelated Bernoulli processes (maximal geometric bias) and the overdamped logarithmic-density dynamics of Cramér, consistent with sieve-induced correlations that partially restore short-gap clustering.

\textbf{Open question:} Quantitative derivation of $R_{\text{Prime}} - R_{\text{Cramér}} \approx 0.22$ from Hardy-Littlewood $k$-tuple correlations would establish the precise contribution of arithmetic structure and close the causal chain from geometric gaps to prime-specific correlations.

Our results reframe the interpretation from ``mysterious prime property'' to ``universal sparse-set phenomenon modulated by arithmetic structure,'' opening new avenues for understanding local prime gap patterns through the lens of stochastic point processes.

\begin{thebibliography}{99}

\bibitem{Rubinstein1994}
M.~Rubinstein and P.~Sarnak, \emph{Chebyshev's bias}, Experimental Mathematics \textbf{3} (1994), 173--197.

\bibitem{Granville2006}
A.~Granville and G.~Martin, \emph{Prime number races}, American Mathematical Monthly \textbf{113} (2006), 1--33.

\bibitem{Cramer1936}
H.~Cramér, \emph{On the order of magnitude of the difference between consecutive prime numbers}, Acta Arithmetica \textbf{2} (1936), 23--46.

\bibitem{HardyLittlewood1923}
G.~H.~Hardy and J.~E.~Littlewood, \emph{Some problems of 'Partitio numerorum'; III: On the expression of a number as a sum of primes}, Acta Mathematica \textbf{44} (1923), 1--70.

\bibitem{Maier1985}
H.~Maier, \emph{Primes in short intervals}, Michigan Math.~J.~\textbf{32} (1985), 221--225.

\end{thebibliography}

\end{document}
